\section{Results}
Across the reported experiments, the intervention improved macro-F1 by 3.2 points.

\section{Discussion}

\subsection{Case A: Unsupported mechanism stack}
The gain may arise because the auxiliary loss smooths the optimization landscape, because
the encoder learns more stable features, or because the decoder becomes better calibrated.
The present comparison did not measure any of these mechanisms and cannot verify which, if
any, produced the gain.

\subsection{Case B: Single hedged explanation with local evidence}
The lower error may reflect reduced prediction variance: across ten seeds, variance fell by
18\%, and removing the ensemble restored the original variance.

\subsection{Case C: Multiple hypotheses with evidence and a discriminator}
Class-balanced augmentation increased rare-class recall by 6.1 points without changing
calibration, whereas temperature scaling reduced expected calibration error from 0.082 to
0.041 without changing recall. Separate ablations distinguished their contributions.

\subsection{Case D: Mechanism explicitly undetermined}
The current design does not identify the mechanism responsible for the 3.2-point difference.

\subsection{Case E: Strong claim earned by a controlled experiment}
In the randomized ablation, disabling retrieval removed the full 3.2-point gain while all
other components and seeds were held fixed. Within this benchmark, retrieval caused the
observed improvement.
